\section{可观测性与运营闭环：tracing、insights、expert answers}

当 Agent 开始高并发服务客户，最贵的问题不是“模型一次答错”，而是“我们不知道它为什么错、在哪个链路错、是否正在批量错”。讲者给出三个关键能力：trace 可视化、对话洞察（insights）、从人工中学习（expert answers）。

\subsection{Tracing：把“黑箱对话”变成“可调系统”}

trace 的价值是把每次动作拆成可解释节点：检索命中、工具调用、策略判定、模型输出、后处理。这样团队可以定位是 retrieval 漏召回、tool schema 不匹配，还是策略守卫误拦截。

\begin{knowledgebox}{可观测性最小字段建议}
每轮记录至少包含：会话 ID、策略版本、提示词版本、工具调用参数与返回、关键时延、输出过滤命中、人工接管标记、最终业务动作状态。没有这些字段，就无法高效做回归与归因。
\end{knowledgebox}

\subsection{Insights 与 Expert Answers：从会话中持续生长}

讲者提到可以对全量对话做开放式提问，例如“升级最多的 5 类原因是什么”“某新品类用户最困惑的问题是什么”。这让客服数据从成本中心转为产品迭代资产。
Expert Answers 的思路是：在人工升级会话中抽取高质量专家处理方式，合成为可复用知识，再回灌给 Agent，降低同类问题再次升级概率。

\begin{importantbox}{运营闭环范式}
生产环境中，最有效的改进路径不是无限调 prompt，而是“观测问题分布 $\rightarrow$ 定位根因 $\rightarrow$ 回灌知识/策略 $\rightarrow$ 再评测验证”。
\end{importantbox}

\subsection{本章小结}

没有可观测性，Agent 只能“凭感觉迭代”；有了 tracing 与 insights，团队才能形成数据驱动的工程闭环，并把人工经验系统化沉淀为可扩展能力。

